\chapter{2024:David Baker and Demis Hassabis and John Jumper}

\label{chap:chemistry2024}

The Nobel Prize in Chemistry for the year 2024 was awarded jointly to David Baker and Demis Hassabis and John Jumper.

\textbf{Motivation:}

for computational protein design (one half, David Baker) and for protein structure prediction (the other half, jointly, Demis Hassabis and John Jumper)

\section{David Baker}

\subsection*{Early Life and Education}

David Baker was born in 1962 in Seattle, Washington, USA.

\subsection*{Career and Major Research}

Baker received his doctorate from the University of California, Berkeley, in 1989. He became a professor of biochemistry at the University of Washington, an investigator of the Howard Hughes Medical Institute, and director of the Institute for Protein Design. He led the development of the Rosetta software for modeling and designing proteins, and in 2003 his group designed a novel protein fold with atomic-level accuracy.

\subsection*{Nobel Prize-winning Work}

The work for which David Baker was awarded the Nobel Prize in Chemistry in 2024 was described by the Nobel Committee as: "for computational protein design".

Baker's de novo design methods allowed proteins with new sequences, folds, and functions to be created from first principles, starting only from the rules of protein structure.

\subsection*{Significance and Impact}

His computational protein design makes it possible to create proteins on demand for medicine, biotechnology, and materials, including designed vaccines and nanomachines.

\subsection*{Later Career and Honors}

Baker continues to direct research in protein design at the University of Washington and leads spin-off companies that apply designed proteins.

\subsection*{Legacy}

Rosetta and the Institute for Protein Design have made computational protein design a practical route to custom biomolecules.

\subsection*{Selected Publications}

"Design of a Novel Globular Protein Fold with Atomic-Level Accuracy" (Science, 2003)

\clearpage

\section{Demis Hassabis}

\subsection*{Early Life and Education}

Demis Hassabis was born in 1976 in London, England.

\subsection*{Career and Major Research}

Hassabis studied computer science at the University of Cambridge and earned a doctorate in cognitive neuroscience from University College London in 2009. He co-founded the artificial intelligence company DeepMind in 2010, where he led the development of AlphaFold, a system that predicts protein structures from amino acid sequences.

\subsection*{Nobel Prize-winning Work}

The work for which Demis Hassabis was awarded the Nobel Prize in Chemistry in 2024 was described by the Nobel Committee as: "for computational protein design (one half, David Baker) and for protein structure prediction (the other half, jointly, Demis Hassabis and John Jumper)".

AlphaFold achieved near-experimental accuracy in predicting protein structures, and in 2020 it solved in days what experimental methods required years to determine.

\subsection*{Significance and Impact}

AlphaFold made reliable structure prediction available for millions of proteins, transforming structural biology and accelerating research in medicine and biology.

\subsection*{Later Career and Honors}

Hassabis is chief executive of Google DeepMind, where he continues to lead research in artificial intelligence.

\subsection*{Legacy}

AlphaFold ended a fifty-year problem in biology and demonstrated the power of machine learning for scientific discovery.

\subsection*{Selected Publications}

"Highly Accurate Protein Structure Prediction with AlphaFold" (Nature, 2021)

\clearpage

\section{John Jumper}

\subsection*{Early Life and Education}

John Jumper was born in 1985 in Little Rock, Arkansas, USA.

\subsection*{Career and Major Research}

Jumper earned a bachelor's degree from Vanderbilt University in 2007 and a doctorate in theoretical chemistry from the University of Chicago in 2017. He joined the artificial intelligence company DeepMind, where he designed the AlphaFold2 system that predicts protein structures with atomic-level accuracy.

\subsection*{Nobel Prize-winning Work}

The work for which John Jumper was awarded the Nobel Prize in Chemistry in 2024 was described by the Nobel Committee as: "for computational protein design (one half, David Baker) and for protein structure prediction (the other half, jointly, Demis Hassabis and John Jumper)".

Jumper's architecture for AlphaFold2 combined machine learning with protein structure geometry, producing predictions that rival experimental structures.

\subsection*{Significance and Impact}

AlphaFold2 provided accurate structures for nearly all proteins of known organisms, a resource that has accelerated research across biology.

\subsection*{Later Career and Honors}

Jumper continues to work on machine learning for biology at Google DeepMind.

\subsection*{Legacy}

The AlphaFold2 method he designed is now a standard tool in structural biology and is used worldwide.

\subsection*{Selected Publications}

"Highly Accurate Protein Structure Prediction with AlphaFold" (Nature, 2021)

\clearpage

\section*{Chapter Summary}

The 2024 prize recognized two advances in protein science: David Baker received one half for computational protein design, and Demis Hassabis and John Jumper shared the other half for the protein structure prediction achieved by AlphaFold. Together these methods let scientists create new proteins and predict the structures of existing ones with unprecedented accuracy.
